%%
%% MScheme Internals: A Technical Report on the Interpreter Implementation
%%
\documentclass[12pt]{report}
\usepackage[margin=1in]{geometry}
\usepackage{longtable}
\usepackage{booktabs}
\usepackage{listings}
\usepackage{xcolor}
\usepackage{hyperref}
\usepackage{textcomp}
\usepackage{amssymb}

\lstdefinelanguage{Scheme}{
  morekeywords={define,lambda,if,cond,else,let,let*,letrec,begin,
    set!,quote,quasiquote,unquote,unquote-splicing,and,or,not,
    case,do,delay,force,map,for-each,apply,call-with-current-continuation,
    call/cc,require-modules,load,macro,unwind-protect,eval,
    current-environment},
  sensitive=true,
  morecomment=[l]{;},
  morestring=[b]",
}

\lstdefinelanguage{Modula3}{
  morekeywords={MODULE,INTERFACE,IMPORT,FROM,PROCEDURE,BEGIN,END,
    VAR,CONST,TYPE,RECORD,OBJECT,METHODS,OVERRIDES,RETURN,
    IF,THEN,ELSE,ELSIF,WHILE,DO,FOR,TO,BY,LOOP,EXIT,REPEAT,UNTIL,
    CASE,OF,WITH,TRY,EXCEPT,FINALLY,RAISE,RAISES,
    NEW,ARRAY,REF,REFANY,SET,BOOLEAN,INTEGER,REAL,LONGREAL,CHAR,TEXT,
    TRUE,FALSE,NIL,BRANDED,REVEAL,EVAL,NOT,AND,OR,IN,
    ISTYPE,NARROW,TYPECODE},
  sensitive=true,
  morecomment=[n]{(*}{*)},
  morestring=[b]",
}

\lstset{
  basicstyle=\ttfamily\small,
  breaklines=true,
  frame=single,
  backgroundcolor=\color{gray!5},
  keywordstyle=\bfseries\color{blue!70!black},
  commentstyle=\itshape\color{green!50!black},
  stringstyle=\color{red!60!black},
  showstringspaces=false,
  columns=flexible,
  tabsize=2,
  xleftmargin=2em,
  framexleftmargin=1.5em,
}

\title{\bfseries\sffamily\fontsize{24.88}{30}\selectfont
  MScheme Internals\\[0.5em]
  \large A Technical Report on the Interpreter Implementation}
\author{Mika Nystr\"om \\ {\tt mika@alum.mit.edu}}
\date{March 2026}

\begin{document}
\maketitle
\thispagestyle{empty}

\newpage
\tableofcontents
\newpage

%======================================================================
\chapter{Introduction}
%======================================================================

MScheme is a Scheme interpreter written in Modula-3 that implements a
dialect largely in line with the Revised$^4$ Report on the Algorithmic
Language Scheme (R4RS).  It originated as a line-for-line translation
of Peter Norvig's JScheme---a compact Scheme interpreter written in
Java---into Modula-3, with subsequent additions including a binding
optimization pass at closure-creation time, Modula-3 interoperability
via a stub generator, and various system extensions.

This report examines the internal implementation of the interpreter
in detail: the type system, the reader, the evaluator's main loop,
the environment and scoping mechanism, the ``partial compilation''
that resolves variable references at closure-creation time, the
primitive system, and memory management.  The Modula-3 stub-generation
system is mentioned where relevant but is not the focus.

\paragraph{Source organization.}
The core interpreter lives in the CM3 package \texttt{mscheme/src/}.
The principal files are:

\begin{center}
\begin{tabular}{ll}
\toprule
\textbf{File} & \textbf{Role} \\
\midrule
\texttt{Scheme.i3}, \texttt{Scheme.m3} & Interpreter object and main eval loop \\
\texttt{SchemeClass.i3} & Revealed private fields of \texttt{Scheme.T} \\
\texttt{SchemeObject.i3} & The universal value type \\
\texttt{SchemeSymbol.i3/.m3} & Symbol (interned atom) representation \\
\texttt{SchemeEnvironment.i3/.m3} & Environment abstraction and scope chain \\
\texttt{SchemeEnvironmentUnsafe.m3} & Unsynchronized environment implementation \\
\texttt{SchemeEnvironmentInstanceRep.i3} & Revealed environment storage layout \\
\texttt{SchemeEnvironmentBinding.i3/.m3} & First-class binding cells \\
\texttt{SchemeClosure.i3/.m3} & Closures, binding optimization \\
\texttt{SchemeClosureClass.i3} & Revealed closure fields \\
\texttt{SchemePair.i3/.m3} & Cons cells \\
\texttt{SchemePrimitive.i3/.m3} & Built-in primitive procedures \\
\texttt{SchemePrimitives.i3} & Standard library macros (Scheme source) \\
\texttt{SchemeInputPort.i3/.m3} & Reader (lexer and parser) \\
\texttt{SchemeProcedure.i3/.m3} & Abstract procedure type \\
\bottomrule
\end{tabular}
\end{center}

%======================================================================
\chapter{The Type System}
\label{ch:types}
%======================================================================

\section{The Universal Value Type}

The foundational design decision is stated in \texttt{SchemeObject.i3}:

\begin{lstlisting}[language=Modula3]
INTERFACE SchemeObject;
TYPE T = REFANY;
END SchemeObject.
\end{lstlisting}

Every Scheme value is a Modula-3 \texttt{REFANY}---a traced, garbage-collected
reference to any heap-allocated object.  This means that any Modula-3 heap
object is already a valid Scheme value without wrapping or marshaling.
Type dispatch throughout the interpreter uses Modula-3's
\texttt{ISTYPE}/\texttt{NARROW} runtime type tests, which check the
object's typecode (a small integer stored in the heap header).

\section{Concrete Types}

The mapping between Scheme types and Modula-3 representations is:

\begin{center}
\begin{longtable}{lll}
\toprule
\textbf{Scheme type} & \textbf{Modula-3 type} & \textbf{Notes} \\
\midrule
\endhead
pair & \texttt{SchemePair.T = REF RECORD first, rest : REFANY END}
     & Cons cell \\
empty list & \texttt{NIL}
           & The Modula-3 null pointer \\
boolean & \texttt{SchemeBoolean.T}
        & Interned singletons \\
fixnum & \texttt{SchemeInt.T = REF INTEGER}
       & Cached $[-256..256]$ \\
bignum & \texttt{Mpz.T}
       & GMP arbitrary-precision \\
rational & \texttt{SchemeRational.T}
         & Exact num/den as \texttt{Mpz.T} \\
inexact real & \texttt{SchemeLongReal.T = REF LONGREAL}
             & IEEE~754 double \\
mpfr float & \texttt{SchemeMpfr.T}
           & MPFR, configurable precision \\
complex & \texttt{SchemeComplex.T}
        & re/im as \texttt{SchemeObject.T} \\
string & \texttt{SchemeString.T = REF ARRAY OF CHAR}
       & Mutable character array \\
character & \texttt{SchemeChar.T = REF CHAR}
          & Boxed character \\
symbol & \texttt{SchemeSymbol.T = Atom.T}
       & Interned text (see \S\ref{sec:symbols}) \\
vector & \texttt{REF ARRAY OF REFANY}
       & \\
closure & \texttt{SchemeClosure.T}
        & See Chapter~\ref{ch:closures} \\
macro & \texttt{SchemeMacro.T <: SchemeClosure.T}
      & \\
primitive & \texttt{SchemePrimitive.T}
          & See Chapter~\ref{ch:primitives} \\
input port & \texttt{SchemeInputPort.T}
           & Wraps \texttt{Rd.T} \\
output port & \texttt{Wr.T}
            & Directly; no wrapper \\
environment & \texttt{SchemeEnvironment.T}
            & Extends \texttt{NetObj.T} \\
\bottomrule
\end{longtable}
\end{center}

Several points deserve attention:

\begin{itemize}
\item \textbf{Booleans} are interned singletons accessed via
  \texttt{SchemeBoolean.True()} and \texttt{SchemeBoolean.False()}.
  Truth testing is \texttt{x $\neq$ False()}---anything other than
  \texttt{\#f} is truthy, per the Scheme specification.

\item \textbf{Numbers} use a full numeric tower: exact integers
  (\texttt{SchemeInt.T} for fixnums, \texttt{Mpz.T} for bignums),
  exact rationals (\texttt{SchemeRational.T}), inexact reals
  (\texttt{SchemeLongReal.T}), MPFR floats (\texttt{SchemeMpfr.T}),
  and complex numbers (\texttt{SchemeComplex.T}).  The
  \texttt{SchemeNumber} module provides unified arithmetic dispatch.

\item \textbf{Output ports} are raw Modula-3 \texttt{Wr.T} writers.
  A Modula-3 writer \emph{is} a Scheme output port with no wrapping.

\item \textbf{Environments} extend \texttt{NetObj.T}, meaning they can in
  principle be shared across address spaces via Modula-3's network objects
  protocol.
\end{itemize}

\section{Symbols and Interning}
\label{sec:symbols}

Symbols are represented as CM3 atoms:

\begin{lstlisting}[language=Modula3]
(* SchemeSymbol.i3 *)
TYPE T = Atom.T;
CONST Symbol = Atom.FromText;
CONST ToText = Atom.ToText;
\end{lstlisting}

The \texttt{Atom} module in the CM3 standard library maintains a
process-wide hash table of interned strings.  \texttt{Atom.FromText(t)}
returns the unique, canonical \texttt{Atom.T} for the text \texttt{t};
calling it twice with equal text returns the \emph{same pointer}.
Consequently:

\begin{enumerate}
\item Symbol identity is \textbf{pointer equality} (\texttt{=} in
  Modula-3).  No string comparison is ever needed after interning.
\item Symbols created by different reader invocations, in different
  threads, or by generated stubs all resolve to the same object.
\item The special-form symbols used by the evaluator (\texttt{SYMquote},
  \texttt{SYMif}, \texttt{SYMlambda}, etc.) are interned once at module
  initialization time and thereafter compared by pointer identity in the
  main eval loop.
\end{enumerate}

The helper \texttt{SymEq(a, b)} interns \texttt{b} and does pointer
comparison---this is used in display code where a text literal is more
readable than a pre-interned variable.

%======================================================================
\chapter{The Reader}
\label{ch:reader}
%======================================================================

The reader is implemented in \texttt{SchemeInputPort.m3}.  It wraps a
Modula-3 \texttt{Rd.T} (reader stream) and provides the Scheme
\texttt{read} operation, producing S-expressions from character input.

\section{Internal State}

\begin{lstlisting}[language=Modula3]
REVEAL T = ... OBJECT
    rd : Rd.T;
    isPushedToken, isPushedChar := FALSE;
    pushedToken : Object := NIL;
    pushedChar : INTEGER := -1;
    lNo := 1;             (* line number for error messages *)
    helpfulText : TEXT;   (* e.g., filename, for diagnostics *)
  END;
\end{lstlisting}

The reader maintains a one-token and one-character pushback buffer.
Character input uses \texttt{UnsafeRd.FastGetChar} for performance
(bypasses per-character locking), with a line counter incremented on
each newline.  EOF is represented as the integer $-1$.

\section{Tokenization}

The tokenizer (\texttt{NextToken}) classifies input characters as follows:

\begin{enumerate}
\item \textbf{Whitespace} is skipped.
\item \textbf{Single-character tokens}: \texttt{(}, \texttt{)}, \texttt{'},
  and \texttt{`} return pre-interned sentinel symbols
  (\texttt{LP}, \texttt{RP}, \texttt{SQ}, \texttt{BQ}).
\item \textbf{Comma}: reads one character ahead to distinguish \texttt{,}
  from \texttt{,@}.
\item \textbf{Semicolon}: line comment; skips to end of line, then recurses.
\item \textbf{Double quote}: reads a string literal, handling backslash
  escapes (\texttt{\textbackslash n}, \texttt{\textbackslash t},
  \texttt{\textbackslash\textbackslash}, \texttt{\textbackslash "}),
  and returns a \texttt{SchemeString.T} (\texttt{REF ARRAY OF CHAR}).
\item \textbf{Hash prefix} (\texttt{\#}): dispatches on the following
  character---\texttt{\#t}/\texttt{\#f} for booleans,
  \texttt{\#(} for vectors, \texttt{\#\textbackslash} for character
  literals (with special handling of \texttt{\#\textbackslash space}
  and \texttt{\#\textbackslash newline}).
\item \textbf{Default case} (identifiers and numbers): characters are
  accumulated in a lightweight \texttt{Wx} buffer (a growable
  \texttt{REF ARRAY OF CHAR} with amortized-constant append) until a
  delimiter or whitespace is encountered.  The resulting text is then:
  \begin{itemize}
  \item Parsed as a \texttt{LONGREAL} if it begins with a digit, sign,
    or decimal point, and contains no non-numeric characters other
    than~\texttt{e}.
  \item Otherwise, interned as a symbol via \texttt{Atom.FromText}.
  \end{itemize}
\end{enumerate}

\section{List Reading}

The procedure \texttt{ReadTail2} reads the elements of a list after the
opening parenthesis.  It is iterative (not recursive) to avoid stack
overflow on long lists:

\begin{enumerate}
\item Read tokens in a loop, pushing each back and calling \texttt{DoRead}
  to parse a complete sub-expression.
\item Prepend each element to an accumulator list \texttt{res} (building
  the list in reverse).
\item On encountering \texttt{)}, stop and set the tail to \texttt{NIL}.
\item On encountering \texttt{.}, read one more expression as the dotted
  tail, then expect~\texttt{)}.
\item Finally, destructively reverse the accumulated list with
  \texttt{ReverseD} and splice the dotted tail (if any) onto the
  last~\texttt{rest} field.
\end{enumerate}

An older recursive \texttt{ReadTail} exists in the source but is marked
\texttt{<*UNUSED*>}.

\section{The Top-Level Read Dispatch}

The \texttt{DoRead} procedure dispatches on the first token:

\begin{lstlisting}[language=Modula3]
WITH token = t.nextToken(...) DO
  IF    token = LP    THEN RETURN t.readTail(...)
  ELSIF token = RP    THEN t.warn("Extra ) ignored"); ...
  ELSIF token = SQ    THEN RETURN List2(quote, t.doRead(...))
  ELSIF token = BQ    THEN RETURN List2(quasiquote, t.doRead(...))
  ELSIF token = COM   THEN RETURN List2(unquote, t.doRead(...))
  ELSIF token = COMAT THEN RETURN List2(unquote-splicing, t.doRead(...))
  ELSE  RETURN token
  END
END
\end{lstlisting}

Quote abbreviations (\texttt{'}, \texttt{`}, \texttt{,}, \texttt{,@})
are expanded by the reader into their \texttt{(quote~\ldots)} /
\texttt{(quasiquote~\ldots)} forms.

%======================================================================
\chapter{The Evaluator}
\label{ch:eval}
%======================================================================

The evaluator is the heart of the interpreter.  It lives in
\texttt{Scheme.m3} and consists of two procedures: a thin wrapper
\texttt{Eval} that provides traceback generation, and the core
\texttt{EvalInternal} that implements the actual dispatch loop.

\section{Traceback Generation: \texttt{Eval}}

\texttt{Eval} calls \texttt{EvalInternal} inside a \texttt{TRY}/\texttt{EXCEPT}
block.  If \texttt{EvalInternal} raises the exception \texttt{E(txt)},
\texttt{Eval} appends a traceback line:

\begin{lstlisting}[language=Modula3]
RAISE E(txt & "\n(called from) eval " &
        TruncateText(Stringify(x), MaxPerLine))
\end{lstlisting}

As each stack frame unwinds, the exception text accumulates a trace.
The trace is capped at \texttt{MaxTraceback} (4096 bytes) to prevent
runaway growth.  A special sentinel string \texttt{"CallCC123"} is passed
through without annotation---this is the internal protocol used by
\texttt{call/cc} (see \S\ref{sec:callcc}).

On the first error, \texttt{Eval} saves the faulting environment and
expression in \texttt{t.errorEnvironment} and \texttt{t.errorEvalX} for
post-mortem inspection from the REPL.

Tracebacks can be disabled globally by setting the environment variable
\texttt{NOMSCHEMETRACEBACKS}.

\section{The Main Loop: \texttt{EvalInternal}}
\label{sec:evalinternal}

\texttt{EvalInternal} takes three arguments: the interpreter \texttt{t},
an expression \texttt{x}, and an environment \texttt{env}.  Its overall
structure is an infinite \texttt{LOOP}:

\begin{lstlisting}[language=Modula3]
PROCEDURE EvalInternal(t : T; x : Object;
                       env : SchemeEnvironment.Instance) : Object
  RAISES { E } =
VAR
  envIsLocallyMade := FALSE;
  savedEnv : SchemeEnvironment.Instance := NIL;
BEGIN
  LOOP
    (* dispatch on the type of x *)
    ...
  END
END EvalInternal;
\end{lstlisting}

The \texttt{LOOP} is the mechanism for tail-call optimization: instead of
making a recursive call in tail position, the code assigns new values to
\texttt{x} and \texttt{env} and goes around the loop again.  The local
variables \texttt{envIsLocallyMade} and \texttt{savedEnv} support
environment recycling across tail calls (see \S\ref{sec:envrecycle}).

\subsection{Type Dispatch}

At the top of each iteration, the evaluator dispatches on the type
of~\texttt{x}:

\begin{lstlisting}[language=Modula3]
IF x = NIL THEN
  RETURN NIL
ELSIF ISTYPE(x, Symbol) THEN
  RETURN env.lookup(x)
ELSIF ISTYPE(x, Binding) THEN
  RETURN NARROW(x, Binding).get()
ELSIF NOT ISTYPE(x, Pair) THEN
  RETURN x                     (* self-evaluating *)
ELSE
  (* x is a Pair: special form or application *)
  ...
END
\end{lstlisting}

\begin{description}
\item[\texttt{NIL}] The empty list is self-evaluating.
\item[\texttt{Symbol}] A variable reference---looked up via the
  environment's scope chain (\texttt{env.lookup}).
\item[\texttt{Binding}] A pre-resolved variable reference (see
  Chapter~\ref{ch:closures}).  The value is fetched directly from the
  binding cell, bypassing the scope chain entirely.
\item[\texttt{Not a Pair}] Numbers, strings, booleans, characters,
  vectors, and any other non-pair, non-symbol objects are self-evaluating.
\item[\texttt{Pair}] An application or special form.
\end{description}

The \texttt{Binding} case is central to understanding MScheme's
performance.  When a closure is created, free variable references in its
body are replaced with \texttt{SchemeEnvironmentBinding.T} objects that
hold a direct pointer to the environment cell.  At eval time, the
evaluator recognizes these objects and calls \texttt{.get()} in constant
time---no scope-chain traversal is needed.

\subsection{Special Forms}

When \texttt{x} is a pair, the evaluator extracts the first element
(\texttt{fn}) and checks it against pre-interned special-form symbols.
Since symbols are interned atoms, this check is pointer comparison.
The special forms are:

\begin{description}
\item[\texttt{quote}] Returns the second element of the form unevaluated.

\item[\texttt{begin}] Evaluates all sub-expressions except the last,
  then sets \texttt{x} to the last expression and loops (tail-call
  optimization):
\begin{lstlisting}[language=Modula3]
WHILE Rest(args) # NIL DO
  EVAL t.eval(First(args), env);
  args := Rest(args)
END;
x := First(args)   (* tail position *)
\end{lstlisting}

\item[\texttt{define}] Two syntactic forms are supported:
  \texttt{(define var val)} stores the value of \texttt{val} in the
  current environment under \texttt{var}, and
  \texttt{(define (f params...) body...)} is sugar for
  \texttt{(define f (lambda (params...) body...))}.  In both cases,
  if the value is an anonymous procedure, it is named after the
  variable.

\item[\texttt{set!}] Evaluates the value, then walks the scope chain
  to find the frame where the variable is defined and mutates it there.
  Raises an error if the variable is unbound.

\item[\texttt{if}] Evaluates the test expression; then sets \texttt{x}
  to either the consequent or the alternative and loops:
\begin{lstlisting}[language=Modula3]
IF TruthO(t.eval(First(args), env)) THEN
  x := Second(args)
ELSE
  x := Third(args)
END
\end{lstlisting}

\item[\texttt{cond}] Delegates to a helper \texttt{ReduceCond} that
  desugars into an expression the main loop can handle (a \texttt{quote},
  a function application, or a \texttt{begin} form).  The result is
  assigned back to~\texttt{x} for tail-call optimization.  The
  \texttt{=>} syntax is supported: \texttt{(cond (test => proc))}
  applies \texttt{proc} to the test value.

\item[\texttt{lambda}] Creates a new closure (see
  Chapter~\ref{ch:closures}) and marks the current environment as
  ``assigned'' (\texttt{env.assigned := TRUE}), preventing it from
  being recycled by the tail-call optimizer.

\item[\texttt{macro}] Like \texttt{lambda} but creates a
  \texttt{SchemeMacro.T}.  Also sets \texttt{env.assigned := TRUE}.

\item[\texttt{eval}] Double-evaluation: evaluates the argument to
  obtain an expression, then sets \texttt{x} to that expression and
  loops.

\item[\texttt{current-environment}] Returns the current environment
  object as a first-class value.

\item[\texttt{unwind-protect}] A three-argument form
  \texttt{(unwind-protect body cleanup handler)} that evaluates
  \texttt{body} in a \texttt{TRY} block, executes \texttt{cleanup}
  unconditionally (as a \texttt{FINALLY}), and optionally catches
  exceptions via \texttt{handler}.
\end{description}

\subsection{Procedure Application}

If \texttt{fn} is not a special-form symbol, it is evaluated
to obtain a procedure.  The evaluator then uses a \texttt{TYPECASE}
to dispatch on the procedure type.

\subsubsection{Closure Application and Tail-Call Optimization}
\label{sec:closure-apply}

When a closure is called, the evaluator avoids a recursive call by
setting \texttt{x~:= c.body} and \texttt{env} to a new environment
frame binding the closure's parameters to the (evaluated) arguments,
then continuing the loop.  This is the core tail-call optimization
for closures.

The environment for the call is managed carefully to minimize allocation:

\begin{lstlisting}[language=Modula3]
VAR newEnv : SchemeEnvironment.Instance;
BEGIN
  IF savedEnv # NIL THEN
    newEnv := savedEnv;  savedEnv := NIL
  ELSE
    newEnv := NEW(SchemeEnvironment.Unsafe)
  END;

  IF envIsLocallyMade AND NOT env.assigned THEN
    savedEnv := env       (* save for reuse on next tail call *)
  END;

  env := newEnv.initEval(c.params, args, env, t, c.env);
  envIsLocallyMade := TRUE;

  IF savedEnv # NIL AND savedEnv.assigned THEN
    savedEnv := NIL       (* captured; cannot reuse *)
  END
END
\end{lstlisting}

The variable \texttt{savedEnv} caches a previously allocated environment
frame.  If one is available, it is reused rather than allocating a new
object.  The current environment can be saved for future reuse only if
(a)~it was created inside this invocation of \texttt{EvalInternal}
(\texttt{envIsLocallyMade}) and (b)~it has not been captured by a closure
(\texttt{NOT env.assigned}).  This recycling is safe because a locally-made,
unassigned environment has no external references---it was created for
the previous tail call and is now dead.

After \texttt{initEval}, the code re-checks \texttt{savedEnv.assigned}
because \texttt{initEval} evaluates the arguments, which might create
closures that capture the saved environment.

\subsubsection{Macro Expansion}

When the evaluated operator is a macro, the evaluator calls
\texttt{m.expand(t, x, args)}, which applies the macro transformer
to the arguments \emph{unevaluated}.  The returned expression is
assigned to~\texttt{x} and the loop continues, giving tail-call
optimization for the expansion result.

\subsubsection{Non-Closure Procedure Application}

For primitives and other procedure types, the evaluator provides
arity-specific fast paths:

\begin{lstlisting}[language=Modula3]
TYPECASE args OF
  NULL => (* zero args: general path *)
| Pair(pp) =>
    IF pp.rest = NIL THEN
      RETURN p.apply1(t, EvalInternal(t, pp.first, env))
    ELSIF ISTYPE(pp.rest, Pair) AND
          NARROW(pp.rest, Pair).rest = NIL THEN
      RETURN p.apply2(t,
        EvalInternal(t, pp.first, env),
        EvalInternal(t, NARROW(pp.rest, Pair).first, env))
    END
ELSE END;
RETURN p.apply(t, t.evalList(args, env))
\end{lstlisting}

The \texttt{apply1} and \texttt{apply2} methods avoid constructing
an argument list entirely.  For the one-argument and two-argument
cases---which account for a large fraction of primitive
calls---no cons cells are allocated for the argument list.

The general \texttt{apply} path calls \texttt{evalList} to evaluate
all arguments and build a proper list, then passes it to
\texttt{p.apply}.

\subsection{Argument List Evaluation: \texttt{EvalList2}}
\label{sec:evallist}

The procedure \texttt{EvalList2} evaluates each element of a list and
builds a new list of results.  It uses \texttt{GetCons} (from the
free-pair pool; see \S\ref{sec:conspools}) for cons cell allocation and
builds the result list tail-first to avoid reversing:

\begin{lstlisting}[language=Modula3]
WITH new = GetCons(t) DO
  new.first := t.eval(pair.first, env);
  new.rest := NIL;
  IF res = NIL THEN
    ptr := new;  res := ptr
  ELSE
    ptr.rest := new;  ptr := new
  END
END
\end{lstlisting}

A pointer \texttt{ptr} tracks the last cons cell, allowing constant-time
append to the tail.

%======================================================================
\chapter{The Environment System}
\label{ch:env}
%======================================================================

\section{Type Hierarchy}

The environment type hierarchy uses Modula-3's multi-level opaque type
revelation:

\begin{center}
\begin{tabular}{ll}
\toprule
\textbf{Type} & \textbf{Adds} \\
\midrule
\texttt{SchemeEnvironmentSuper.T = NetObj.T} & Network object root \\
\texttt{SchemeEnvironment.T} & \texttt{lookup}, \texttt{bind}, \texttt{define}, \texttt{set}, \texttt{defPrim} \\
\texttt{SchemeEnvironment.PubInstance} & \texttt{assigned} flag; \texttt{init}, \texttt{initEval}, \texttt{initEmpty} \\
\texttt{SchemeEnvironmentClass.Private} & \texttt{put}, \texttt{get} (internal storage API) \\
\texttt{SchemeEnvironmentInstanceRep.Rep} & \texttt{dictionary}, \texttt{quick}, \texttt{parent}, \texttt{dead} \\
\texttt{SchemeEnvironment.Instance} & The fully-revealed concrete type \\
\bottomrule
\end{tabular}
\end{center}

Two concrete subtypes exist: \texttt{Unsafe} (no synchronization, used
for local frames and the default global environment) and \texttt{Safe}
(mutex-protected, for shared or network-exported environments).

\section{Storage: The Quick Array and Hash Table}
\label{sec:quickarray}

Each environment frame stores bindings in a two-tier structure, defined
in \texttt{SchemeEnvironmentInstanceRep.i3}:

\begin{lstlisting}[language=Modula3]
CONST QuickVars = 10;
TYPE QuickMap = RECORD var : Symbol; val : Object END;

TYPE Rep = ... OBJECT
    dictionary : AtomRefTbl.T;
    quick : ARRAY [0..QuickVars - 1] OF QuickMap;
    parent : SchemeEnvironment.T;
    dead := FALSE;
    assigned := FALSE;
  END;
\end{lstlisting}

The \texttt{quick} array holds the first 10 bindings inline as
\texttt{(symbol, value)} pairs.  This is a linear-scanned array
but is cache-friendly and fast for small environments.  Since most
lambda parameter lists have fewer than~10 entries, the hash table
is rarely needed for local frames.

The \texttt{dictionary} field is a hash table (\texttt{AtomRefTbl.T},
keyed by \texttt{Atom.T}) that is lazily allocated only when the quick
array overflows.

\subsection{\texttt{put}: Storing a Binding}

\begin{lstlisting}[language=Modula3]
PROCEDURE UnsafePut(t : Instance; var : Symbol;
                    READONLY val : Object) =
BEGIN
  FOR i := FIRST(t.quick) TO LAST(t.quick) DO
    IF t.quick[i].var = var OR t.quick[i].var = NIL THEN
      t.quick[i].var := var;
      t.quick[i].val := val;
      RETURN
    END
  END;
  (* quick array full: spill to hash table *)
  IF t.dictionary = NIL THEN
    t.dictionary := NEW(AtomRefTbl.Default).init()
  END;
  EVAL t.dictionary.put(var, val)
END UnsafePut;
\end{lstlisting}

The scan looks for either a matching variable (update in place) or a
\texttt{NIL} slot (insert).  The comparison \texttt{t.quick[i].var = var}
is pointer equality on interned atoms.

\subsection{\texttt{get}: Retrieving a Binding}

\begin{lstlisting}[language=Modula3]
PROCEDURE UnsafeGet(t : Instance; var : Symbol;
                    VAR val : Object) : BOOLEAN =
BEGIN
  IF var = NIL THEN RETURN FALSE END;
  FOR i := FIRST(t.quick) TO LAST(t.quick) DO
    IF t.quick[i].var = var THEN
      val := t.quick[i].val; RETURN TRUE
    END
  END;
  IF t.dictionary # NIL THEN
    RETURN t.dictionary.get(var, val)
  ELSE
    RETURN FALSE
  END
END UnsafeGet;
\end{lstlisting}

Linear scan of the quick array, then fallback to the hash table.

\subsection{\texttt{lookup}: Scope Chain Traversal}

\begin{lstlisting}[language=Modula3]
PROCEDURE Lookup(t : Instance; symbol : Symbol) : Object
  RAISES { E } =
VAR o : Object;
BEGIN
  <*ASSERT NOT t.dead*>
  IF t.get(symbol, o) THEN RETURN o END;
  IF t.parent # NIL THEN
    RETURN t.parent.lookup(symbol)
  ELSE
    RETURN Error("Unbound variable in lookup: " &
                 Stringify(symbol))
  END
END Lookup;
\end{lstlisting}

Classic lexical scoping: search the local frame, then walk up the
parent chain.  The \texttt{dead} assertion guards against use-after-recycle
bugs (see \S\ref{sec:envrecycle}).

\section{\texttt{define} and \texttt{set!}}

\texttt{Define} always stores in the \emph{local} frame.  It also
names anonymous procedures:

\begin{lstlisting}[language=Modula3]
PROCEDURE Define(t : Instance; var, val : Object) : Object =
BEGIN
  t.put(var, val);
  TYPECASE val OF SchemeProcedure.T(p) =>
    IF p.name = "anonymous procedure" THEN
      p.name := SchemeSymbol.ToText(var)
    END
  ELSE END;
  ...
END Define;
\end{lstlisting}

\texttt{Set} walks the scope chain to find the frame where the variable
is defined, then mutates it there.  If the variable is unbound everywhere,
it raises an error.

\section{The \texttt{assigned} Flag}
\label{sec:assigned}

The boolean field \texttt{assigned} on each environment frame records
whether the frame has been captured by a closure.  It is set to
\texttt{TRUE} whenever a \texttt{lambda} or \texttt{macro} form is
evaluated in that environment, and also during the closure binding
optimization in \texttt{SchemeClosure.Init}.

The flag serves a single purpose: it tells the tail-call optimizer in
\texttt{EvalInternal} whether an environment frame can be recycled.
An environment that has been ``assigned'' (captured) may be reachable
from a closure and must not be overwritten.

\section{Environment Initialization}
\label{sec:envinit}

Environments are initialized via three methods:

\begin{description}
\item[\texttt{initEmpty(parent)}] Clears all quick slots to \texttt{NIL},
  nils out the dictionary, and sets the parent pointer.  Used when
  recycling an environment frame.

\item[\texttt{init(vars, vals, parent)}] Creates a frame and populates
  it by walking two parallel lists: parameter names and argument values.
  If \texttt{vars} is a bare symbol (rather than a pair), it receives
  all remaining values as a list (rest-argument binding).

\item[\texttt{initEval(vars, argsToEval, evalEnv, interp, parent)}]
  Like \texttt{init}, but evaluates each argument expression in
  \texttt{evalEnv} before binding.  This is the path used by the
  tail-call optimizer, and it avoids constructing an intermediate
  argument list that \texttt{evalList} would produce.
\end{description}

The \texttt{initEval} method calls \texttt{InitDictEval2}, which
evaluates arguments one at a time:

\begin{lstlisting}[language=Modula3]
LOOP
  IF vars = NIL AND argsToEval = NIL THEN RETURN TRUE
  ELSIF vars # NIL AND ISTYPE(vars, Symbol) THEN
    t.put(vars, interp.evalList(argsToEval, evalEnv));
    RETURN TRUE                     (* rest-arg *)
  ELSIF vars # NIL AND argsToEval # NIL AND ... THEN
    t.put(varp.first, interp.eval(argp.first, evalEnv));
    vars := varp.rest;
    argsToEval := argp.rest
  ELSE
    RETURN FALSE                    (* arity mismatch *)
  END
END
\end{lstlisting}

This avoids the cons-cell allocation that \texttt{evalList} would
require: arguments are evaluated and stored directly into the new
environment frame.

\section{Environment Recycling}
\label{sec:envrecycle}

As described in \S\ref{sec:closure-apply}, the tail-call optimizer
maintains a \texttt{savedEnv} variable that caches one environment
frame for reuse.  The frame is reused by calling \texttt{initEmpty}
to clear its contents, rather than allocating a new object.

The \texttt{dead} field (set during \texttt{initEmpty} and cleared
afterward, or used as a debugging assertion) helps catch bugs where
a recycled environment is accidentally accessed.

%======================================================================
\chapter{The Binding Optimization}
\label{ch:closures}
%======================================================================

This chapter describes the ``partial compilation'' that occurs when a
closure is created.  It is the most distinctive optimization in MScheme
and the key to understanding its performance characteristics.

\section{Overview}

When the evaluator encounters a \texttt{lambda} form, it does not
simply store the parameter list and body verbatim.  Instead, it walks
the body and \textbf{replaces free variable references (symbols) with
\texttt{SchemeEnvironmentBinding.T} objects} that point directly to the
environment cell where each variable lives.  At subsequent evaluation
time, these binding objects are recognized by the evaluator's type
dispatch and resolved in constant time---no scope-chain walk is
needed.

This is a form of compile-time name resolution: the ``compilation''
happens once, when the closure is created, and the benefit is
amortized over every subsequent call to the closure.

\section{The Closure Object}

A closure's private fields (from \texttt{SchemeClosureClass.i3}):

\begin{lstlisting}[language=Modula3]
Private = SchemeClosure.Public OBJECT
  params, body : SchemeObject.T;
  env          : SchemeEnvironment.Instance;
END;
\end{lstlisting}

\texttt{params} is the parameter list (a proper list of symbols, or a
list ending with a dotted symbol for rest arguments).  \texttt{body} is
the (possibly transformed) body expression.  \texttt{env} is the
lexical environment at the point of closure creation.

\section{The \texttt{Init} Procedure}

\texttt{SchemeClosure.Init} in \texttt{SchemeClosure.m3} takes a
\texttt{bind} parameter.  When \texttt{bind} is \texttt{TRUE} (which
it is for \texttt{lambda}; macros pass \texttt{FALSE}), the binding
optimization is performed.

The core logic is in the nested procedure \texttt{DoBindings}, which
walks the body recursively:

\begin{lstlisting}[language=Modula3]
PROCEDURE DoBindings(body : Pair;
                     VAR wasDefine : BOOLEAN) : Object =
\end{lstlisting}

For each element of the body:

\begin{enumerate}
\item If it is a \textbf{Pair} (sub-expression): recurse.

\item If it is a \textbf{Symbol in the first (operator) position}:
  \begin{itemize}
  \item If it is a special-form symbol (\texttt{IsSpecialForm(sym)}),
    bail out---return the body unchanged.  Special forms have non-standard
    evaluation rules and their arguments must not be pre-resolved.
  \item If it is a macro or unknown binding (\texttt{IsMacroOrUnknown(sym)}),
    bail out.  Unknown symbols might refer to macros not yet defined;
    pre-resolving them would be incorrect.
  \item If it is \texttt{define}, set \texttt{wasDefine := TRUE} and
    bail out.  A \texttt{define} introduces new bindings that would
    invalidate subsequent pre-resolutions.
  \end{itemize}

\item If it is a \textbf{Symbol that is a parameter of this lambda}
  (\texttt{MemberP(sym, params)}): leave it alone.  It will be bound
  in the new local environment at call time.

\item If it is a \textbf{Symbol that is a free variable}: call
  \texttt{env.bind(sym)} to obtain a \texttt{SchemeEnvironmentBinding.T}
  object, and replace the symbol in the body with this object.

\item If \texttt{env.bind} fails (variable not found in any enclosing
  scope): bail out and return the original body unchanged.

\item \textbf{Non-symbol, non-pair elements} (numbers, strings, etc.)
  are left as-is; they are self-evaluating.
\end{enumerate}

The macro-or-unknown check is:

\begin{lstlisting}[language=Modula3]
PROCEDURE IsMacroOrUnknown(sym : SchemeSymbol.T) : BOOLEAN =
BEGIN
  IF NOT env.haveBinding(sym) THEN
    RETURN TRUE    (* unknown = might be a macro later *)
  ELSE
    RETURN ISTYPE(env.lookup(sym), SchemeMacro.T)
  END
END IsMacroOrUnknown;
\end{lstlisting}

This is conservative: if a symbol is not currently bound, it might be
a macro that will be defined later (e.g., by a \texttt{load}).  In
that case, pre-resolving it would be wrong.  The code errs on the side
of caution and leaves the symbol as-is.

\subsection{Multi-Expression Bodies}

The outer wrapper \texttt{MapBindings} iterates over each top-level
expression in a multi-expression body.  If any expression is or
contains a \texttt{define}, processing stops and all subsequent
expressions are left unchanged.  This is because \texttt{define}
introduces new bindings into the scope, potentially shadowing names
that were already pre-resolved.

\subsection{Body Wrapping}

After binding resolution, the body is packaged:

\begin{lstlisting}[language=Modula3]
env.assigned := TRUE;
t.params := params;
t.env := env;
IF body # NIL AND ISTYPE(body, Pair)
   AND Rest(body) = NIL THEN
  t.body := First(DoBindings(body, dummy))
ELSIF ISTYPE(body, Pair) THEN
  t.body := Cons(SYMbegin, MapBindings(body))
ELSE
  t.body := Cons(SYMbegin, body)
END;
\end{lstlisting}

A single-expression body is unwrapped from its list (avoiding a
\texttt{begin} form).  Multi-expression bodies are wrapped in
\texttt{(begin~\ldots)}.

\section{The Binding Object}
\label{sec:bindingobj}

\texttt{SchemeEnvironmentBinding.T} is an abstract object type:

\begin{lstlisting}[language=Modula3]
TYPE T = OBJECT METHODS
    name()  : SchemeSymbol.T;
    env()   : SchemeObject.T;
    get()   : SchemeObject.T;
    setB(to : SchemeObject.T);
  END;
\end{lstlisting}

The \texttt{Unsafe} environment provides an optimized implementation
called \texttt{MyBinding}:

\begin{lstlisting}[language=Modula3]
TYPE MyBinding = SchemeEnvironmentBinding.T OBJECT
    e : Unsafe;
    s : Symbol;
    q : [-1..LAST(CARDINAL)] := -1;
  OVERRIDES
    get  := SBGet;
    setB := SBSetB;
  END;
\end{lstlisting}

The field \texttt{q} is the index into the quick array at the time
the binding was created.  When \texttt{q $\neq -1$}, the binding's
\texttt{get} method is a single indexed read:

\begin{lstlisting}[language=Modula3]
PROCEDURE SBGet(sb : MyBinding) : Object =
VAR v : Object;
BEGIN
  IF sb.q # -1 THEN
    RETURN sb.e.quick[sb.q].val    (* direct array access *)
  END;
  WITH gotit = sb.e.get(sb.s, v) DO
    <*ASSERT gotit*>
    RETURN v
  END
END SBGet;
\end{lstlisting}

This is the fastest possible variable access in MScheme: one pointer
dereference to the environment object, one indexed read from a
fixed-size array.  No hash lookup, no scope chain traversal.

The \texttt{Safe} environment provides a simpler \texttt{SimpleBinding}
that does not cache the quick-array index but still avoids the scope
chain walk by holding a direct pointer to the environment frame.

\section{How Bindings Are Created}

The \texttt{GetBinding} method on an \texttt{Unsafe} environment:

\begin{lstlisting}[language=Modula3]
PROCEDURE GetBinding(t : Instance; sym : Symbol) : Binding
  RAISES { E } =
BEGIN
  FOR i := FIRST(t.quick) TO LAST(t.quick) DO
    IF t.quick[i].var = sym THEN
      RETURN NEW(MyBinding, e := t, q := i, s := sym)
    END
  END;
  IF t.get(sym, o) THEN
    RETURN NEW(MyBinding, e := t, s := sym)
  END;
  IF t.parent # NIL THEN
    RETURN t.parent.bind(sym)
  ELSE
    RETURN Error("Unbound variable: " & Stringify(sym))
  END
END GetBinding;
\end{lstlisting}

If the variable is found in the quick array, the binding caches the
index \texttt{q}.  If it is in the hash table, \texttt{q} remains
$-1$ and subsequent accesses go through \texttt{get}.  If it is not
in the local frame, the search continues up the parent chain---but
the returned binding will point to whichever ancestor frame holds
the variable.

\section{When Binding Fails}

The binding optimization is deliberately conservative.  It bails out
(leaving the original symbol in place) in several situations:

\begin{itemize}
\item The operator position of a form is a special-form symbol.
\item The operator is a macro or is currently unbound (might become
  a macro later).
\item A \texttt{define} appears in the body, potentially introducing
  new bindings.
\item A free variable is not found in any enclosing scope.
\end{itemize}

In all these cases, the closure falls back to ordinary symbol lookup
at evaluation time.  The optimization is purely opportunistic; its
absence does not affect correctness.

\section{Interaction with \texttt{set!}}

The \texttt{setB} method on a binding mutates the value in the
environment frame.  For \texttt{MyBinding}:

\begin{lstlisting}[language=Modula3]
PROCEDURE SBSetB(sb : MyBinding; to : Object) =
BEGIN
  IF sb.q # -1 THEN
    sb.e.quick[sb.q].val := to
  ELSE
    sb.e.put(sb.s, to)
  END
END SBSetB;
\end{lstlisting}

This means that \texttt{set!} on a pre-resolved variable is also
a constant-time operation---no scope chain walk is needed.

%======================================================================
\chapter{The Primitive System}
\label{ch:primitives}
%======================================================================

\section{The Procedure Type}

All callable objects in MScheme extend \texttt{SchemeProcedure.T}:

\begin{lstlisting}[language=Modula3]
TYPE T <: Public;
Public = OBJECT METHODS
    format() : TEXT;
    apply(interp : Scheme.T; args : Object) : Object
      RAISES { E };
    apply1(interp : Scheme.T; a1 : Object) : Object
      RAISES { E };
    apply2(interp : Scheme.T; a1, a2 : Object) : Object
      RAISES { E };
  END;
\end{lstlisting}

The \texttt{apply1} and \texttt{apply2} methods exist for the
arity-specific optimizations described in \S\ref{sec:evalinternal}.
The default implementations wrap their arguments into a list and
call~\texttt{apply}; \texttt{SchemePrimitive.T} overrides them to
use the cons-cell pool.

Each procedure also carries a \texttt{name} field (defaulting to
\texttt{"anonymous procedure"}) used in error messages and tracebacks.

\section{Primitive Registration}

A \texttt{SchemePrimitive.T} carries an integer ID, a reference to
its \texttt{Definer}, and min/max argument counts.  Primitives are
registered in layered \texttt{InstallPrimitives} procedures:

\begin{center}
\begin{tabular}{ll}
\toprule
\textbf{Layer} & \textbf{Contents} \\
\midrule
\texttt{InstallSandboxPrimitives} & $\sim$170 safe primitives (R4RS core) \\
\texttt{InstallFileIOPrimitives} & File I/O (\texttt{open-input-file}, etc.) \\
\texttt{InstallNorvigPrimitives} & JScheme extensions (\texttt{new}, \texttt{class}, \texttt{exit}) \\
\texttt{InstallDefaultExtendedPrimitives} & MScheme extensions (\texttt{random}, \texttt{system}, etc.) \\
\bottomrule
\end{tabular}
\end{center}

Each primitive is installed via \texttt{env.defPrim(name, id, definer,
minArgs, maxArgs)}, which interns the name as a symbol, creates a
\texttt{SchemePrimitive.T} with the given ID, and defines it in the
environment.

\section{Dispatch}

The dispatch hub is the procedure \texttt{Prims}, which uses a
\texttt{CASE} statement on an enumeration type \texttt{P}:

\begin{lstlisting}[language=Modula3]
CASE VAL(t.id, P) OF
  P.Eq   => RETURN NumCompare(args, '=')
| P.Lt   => RETURN NumCompare(args, '<')
| P.Car  => RETURN PedanticFirst(x)
| P.Cons => RETURN Cons(x, y, interp)
| ...
END
\end{lstlisting}

For user-defined extension primitives (registered via
\texttt{ExtDefiner.addPrim}), the ID is a random integer above
\texttt{ORD(LAST(P))}, and dispatch goes through a hash table lookup:

\begin{lstlisting}[language=Modula3]
IF t.id > ORD(LAST(P)) THEN
  RETURN NARROW(t.definer, ExtDefiner).apply(t, interp, args)
\end{lstlisting}

\section{Arity-Optimized Dispatch}

\texttt{SchemePrimitive.T} overrides \texttt{apply1} and \texttt{apply2}
to use the cons-cell pool:

\begin{lstlisting}[language=Modula3]
PROCEDURE Apply1(t : T; interp : Scheme.T; a1 : Object)
  : Object RAISES { E } =
VAR d1 := GetCons(interp); free := DefFree;
BEGIN
  d1.first := a1;  d1.rest := NIL;
  WITH res = Prims(t, interp, d1, a1, NIL, free) DO
    IF free THEN ReturnCons(interp, d1) END;
    RETURN res
  END
END Apply1;
\end{lstlisting}

The temporary cons cell is allocated from the pool and returned after
the call.  The \texttt{free} flag tracks whether the primitive captured
the argument list (e.g., \texttt{list} or \texttt{append} would set
\texttt{free := FALSE}), in which case the cell must not be recycled.

\section{Standard Library Macros}

The interface \texttt{SchemePrimitives.i3} contains a single constant
\texttt{Code}: a long \texttt{TEXT} of Scheme source code that is
evaluated at interpreter initialization time.  It defines the following
standard forms as macros:

\begin{center}
\begin{tabular}{lp{9cm}}
\toprule
\textbf{Form} & \textbf{Expansion strategy} \\
\midrule
\texttt{or} & Expands to \texttt{(cond (a) (b) (c) ...)} \\
\texttt{and} & Recursive: \texttt{(if a (and b c) \#f)} \\
\texttt{quasiquote} & Recursive expander tracking nesting depth \\
\texttt{let} & \texttt{((lambda (vars...) body...) vals...)}; supports named \texttt{let} via \texttt{set!} \\
\texttt{let*} & Nested \texttt{let}s, one per binding \\
\texttt{letrec} & Binds all to \texttt{\#f}, then \texttt{set!}s all \\
\texttt{case} & \texttt{cond} with \texttt{member} tests \\
\texttt{do} & \texttt{letrec}-based tail-recursive loop \\
\texttt{delay} & Creates a memoizing thunk via \texttt{make-promise} \\
\texttt{time} & Wraps in \texttt{time-call} primitive \\
\bottomrule
\end{tabular}
\end{center}

These are non-hygienic procedural macros (JScheme-style).  R5RS
hygienic macros (\texttt{syntax-rules}) are available by loading the
bundled \texttt{mbe} library.

Also defined are convenience aliases:
\texttt{call/cc} for \texttt{call-with-current-continuation},
\texttt{first}/\texttt{rest} for \texttt{car}/\texttt{cdr},
\texttt{second} for \texttt{cadr}, and \texttt{third} for \texttt{caddr}.

%======================================================================
\chapter{Memory Management}
\label{ch:memory}
%======================================================================

MScheme relies on the Modula-3 traced-heap garbage collector for all
memory management.  However, it supplements the GC with two
application-level optimizations to reduce allocation pressure.

\section{The Cons Cell Free List}
\label{sec:conspools}

The interpreter object maintains a singly-linked free list of recycled
cons cells:

\begin{lstlisting}[language=Modula3]
(* SchemeClass.i3 *)
Private = Scheme.Public OBJECT
  freePairs : SchemePair.T := NIL;
  ...
END;
\end{lstlisting}

\subsection{\texttt{GetCons}: Allocation}

\begin{lstlisting}[language=Modula3]
PROCEDURE GetCons(t : T) : Pair =
BEGIN
  IF t.freePairs # NIL THEN
    VAR res := t.freePairs;
    BEGIN
      t.freePairs := res.rest;
      RETURN res
    END
  END;
  RETURN NEW(Pair)
END GetCons;
\end{lstlisting}

Pops the head of the free list if available; otherwise allocates a new
pair from the GC heap.

\subsection{\texttt{ReturnCons}: Deallocation}

\begin{lstlisting}[language=Modula3]
PROCEDURE ReturnCons(t : T; cons : Pair) =
BEGIN
  IF cons = NIL THEN RETURN END;
  VAR p : Pair := cons;
  BEGIN
    WHILE p.rest # NIL DO
      p.first := SYMrip;    (* poison *)
      p := p.rest
    END;
    p.first := SYMrip;
    p.rest := t.freePairs;  (* splice onto free list *)
    t.freePairs := cons
  END
END ReturnCons;
\end{lstlisting}

The entire chain of cons cells is returned in one operation.  Each
cell's \texttt{first} field is overwritten with the poison symbol
\texttt{SYMrip} (the string \texttt{"####r.i.p.-dead-cons-cell####"}).
This is a debugging aid: if a recycled cell is accidentally read, the
poison value will be conspicuous in output.

\subsection{Where Recycling Occurs}

Cons cells from the pool are used in:

\begin{itemize}
\item \texttt{EvalList2} (argument list evaluation)
\item \texttt{SchemePrimitive.Apply1} and \texttt{Apply2} (temporary
  argument packaging for primitive dispatch)
\item Various utility functions (\texttt{Cons}, \texttt{List2}, etc.)
  that accept an interpreter parameter
\end{itemize}

The pool is per-interpreter-instance and is explicitly not thread-safe.

\section{Scratch Buffers}

The interpreter object also carries reusable scratch buffers:

\begin{lstlisting}[language=Modula3]
wx : Wx.T := NIL;       (* string builder *)
refseq : RefSeq.T := NIL;  (* sequence builder *)
\end{lstlisting}

These are available to primitives for string construction and sequence
building, avoiding repeated heap allocation for temporary buffers.

%======================================================================
\chapter{The REPL}
\label{ch:repl}
%======================================================================

\section{The Read-Eval-Write Loop}

The REPL is implemented in \texttt{Scheme.ReadEvalWriteLoop}:

\begin{lstlisting}[language=Modula3]
PROCEDURE ReadEvalWriteLoop(t : T; int : Interrupter)
  RAISES { Wr.Failure } =
BEGIN
  t.setInterrupter(int);
  t.bind(SchemeSymbol.Symbol("bang-bang"), NIL);
  LOOP
    Wr.PutText(t.output, ">");
    Wr.Flush(t.output);
    TRY
      WITH x = t.input.read() DO
        IF SchemeInputPort.IsEOF(x) THEN RETURN END;
        WITH res = t.evalInGlobalEnv(x) DO
          EVAL SchemeUtils.Write(res, t.output, TRUE);
          EVAL t.globalEnvironment.set(
            SchemeSymbol.Symbol("bang-bang"), res)
        END
      END
    EXCEPT
      E(e) =>
        Wr.PutText(t.output, "EXCEPTION! " & e & "\n")
    END;
    Wr.PutText(t.output, "\n");
    Wr.Flush(t.output)
  END
END ReadEvalWriteLoop;
\end{lstlisting}

Key details:

\begin{itemize}
\item The prompt is \texttt{">"} (no trailing space).
\item The variable \texttt{bang-bang} holds the result of the last
  successful evaluation, analogous to \texttt{\_} or \texttt{!!}
  in other REPLs.
\item Reading is done via \texttt{t.input.read()}, which invokes
  the full Scheme reader described in Chapter~\ref{ch:reader}.
\item Evaluation is always in the global environment.
\item Results are displayed with \texttt{SchemeUtils.Write} in
  ``quoted'' mode (\texttt{TRUE} second argument), meaning strings
  are printed with quotes and special characters escaped.
\item Exceptions are caught and displayed as
  \texttt{"EXCEPTION!~"} followed by the error text, which includes
  accumulated tracebacks.
\item The loop terminates only on EOF from the input port.
\end{itemize}

\section{Interpreter Initialization}

The initialization sequence (\texttt{Init2}) proceeds as follows:

\begin{enumerate}
\item Create the input port (wrapping a \texttt{Rd.T}) and set the
  output writer.
\item Create or accept a global environment (\texttt{Unsafe} by default).
\item Bind the search path variable \texttt{**scheme-load-path**} from the
  environment variable \texttt{MSCHEMEPATH} (default: \texttt{"."}).
  The path is colon-separated.
\item Install primitives via the configured \texttt{Definer}.
\item Evaluate the standard library macros from \texttt{SchemePrimitives.Code}.
\item Load any user-specified files.
\item Bind \texttt{*the-global-environment*} to the global environment
  itself.
\end{enumerate}

\section{File Loading}

\texttt{LoadFile} searches for files in this order:

\begin{enumerate}
\item If the path is absolute (\texttt{/}), open directly.
\item If it starts with \texttt{\~{}}, expand the home directory.
\item Otherwise, search the \texttt{**scheme-load-path**} list, trying
  each directory as a prefix.
\item If not found on disk, try the \texttt{SchemeDefsBundle}---a
  resource bundle of \texttt{.scm} files compiled into the binary
  via \texttt{m3bundle}.
\end{enumerate}

%======================================================================
\chapter{Continuations and \texttt{call/cc}}
\label{sec:callcc}
%======================================================================

MScheme inherits the main limitation of JScheme: because the Modula-3
call stack is used as the Scheme call stack, full
\texttt{call-with-current-continuation} is not available.
The implementation uses an exception-based mechanism.

When \texttt{call/cc} is invoked, a \texttt{SchemeContinuation.T} is
created that, when called, raises the exception \texttt{Scheme.E} with
the sentinel text \texttt{"CallCC123"}.  The \texttt{call/cc} primitive
catches this specific exception and extracts the continuation value.

This means continuations can only be used for \emph{escape}
(unwinding the stack); they cannot be re-invoked after the dynamic
extent of the \texttt{call/cc} has ended.  The sentinel string
\texttt{"CallCC123"} is recognized by the traceback mechanism in
\texttt{Eval} and passed through without annotation.

%======================================================================
\chapter{Summary of Design Patterns}
\label{ch:summary}
%======================================================================

The following patterns characterize MScheme's implementation:

\begin{enumerate}
\item \textbf{Tail-call optimization via loop.}  The main \texttt{LOOP}
  in \texttt{EvalInternal} implements proper tail calls for
  \texttt{begin}, \texttt{if}, \texttt{cond}, closure application,
  and macro expansion by mutating \texttt{x} and \texttt{env} and
  continuing the loop.

\item \textbf{Binding pre-resolution.}  At closure creation time, free
  variable symbols are replaced with binding objects that hold direct
  pointers to the environment slot.  This reduces variable lookup from
  $O(\text{scope-depth} \times 10)$ to~$O(1)$.

\item \textbf{Quick array.}  The first 10 bindings in each environment
  frame are stored in a fixed inline array, avoiding hash-table
  allocation for small environments (which is nearly all of them).

\item \textbf{Environment recycling.}  Locally-created environment
  frames that have not been captured by closures (\texttt{assigned =
  FALSE}) are saved and reused on subsequent tail calls.

\item \textbf{Cons cell free list.}  A per-interpreter free list of
  cons cells reduces GC pressure.  Cells are poisoned on return for
  debugging safety.

\item \textbf{Arity-specific dispatch.}  One-argument and two-argument
  primitive calls use dedicated \texttt{apply1}/\texttt{apply2} paths
  that avoid argument-list construction entirely.

\item \textbf{Symbol interning.}  All symbols are Modula-3 atoms with
  process-wide uniqueness, making symbol comparison a single pointer
  test.

\item \textbf{\texttt{REFANY} as the universal type.}  The identity
  \texttt{SchemeObject.T = REFANY} means Modula-3 objects pass into
  and out of Scheme without marshaling, and the Modula-3 GC manages
  both Scheme and Modula-3 data uniformly.
\end{enumerate}

%======================================================================
\chapter*{References}
\addcontentsline{toc}{chapter}{References}
%======================================================================

\begin{enumerate}
\item William Clinger and Jonathan Rees (eds.),
  \emph{Revised$^4$ Report on the Algorithmic Language Scheme},
  ACM Lisp Pointers~IV(3), July--September 1991.

\item Peter Norvig, ``JScheme: Scheme in Java,''
  \url{https://norvig.com/jscheme.html}, 1998.
  MScheme originated as a line-for-line translation of JScheme into
  Modula-3.  The eval-loop structure, the use of the host language's
  call stack as the Scheme call stack, the procedural macro system,
  the \texttt{call/cc} limitation, and the standard-library macros
  defined as source text are all inherited from JScheme.

\item Peter Norvig, ``(An ((Even Better) Lisp) Interpreter (in Python)),''
  \url{https://norvig.com/lispy2.html}, 2010.
  Describes a similar small-Scheme-in-a-host-language design,
  including the trampoline-based tail-call optimization pattern
  that MScheme also uses.

\item Greg Nelson (ed.), \emph{Systems Programming with Modula-3},
  Prentice Hall, 1991.
  The definitive reference for the Modula-3 language and its design
  rationale, including the type system, opaque types, garbage
  collection, and the exception model.

\item Samuel P.\ Harbison, \emph{Modula-3}, Prentice Hall, 1992.
  A practical reference for Modula-3 language features used throughout:
  \texttt{REFANY}, traced references, \texttt{BRANDED}, object types
  with method overrides, and multi-level type revelation.
\end{enumerate}

\end{document}
